\section{Vocabulaire}

\begin{definitionbox}
\begin{itemize}[label=\textbullet]
    \item On appelle \bb{population statistique} l'ensemble 
        des \bb{individus} sur lesquels porte l'étude statistique. 
    \item On appelle \bb{effectif total} de la population statistique le
        nombre d'individus de cette population. 
    \item On appelle \bb{variable statistique} ou \bb{caractère}, la chose que l'on étudie
        et qui est commune à tous les individus de la population. L'ensemble des
        résultats s'appelle série statistique.
    \item Une \bb{variable statistique} peut être :
        \begin{itemize}[label=\textbullet]
            \item \answer{qualitative} si les valeurs prises sont des catégories
                ou des qualités.
            \item \answer{quantitative} si les valeurs prises sont des nombres.
        \end{itemize}
    \item On appelle \bb{effectif} associé à une valeur de la variable, le nombre de
        fois où cette valeur apparait dans la série.
\end{itemize}
\end{definitionbox}

\begin{exemplebox}
    Pour se rendre au collège

    \begin{itemize}[label=\textbullet]
        \item 50 élèves utilisent un deux roues,
        \item 250 élèves utilisent les transports en commun,
        \item 125 élèves se déplacent à pieds,
        \item 75 élèves sont déposés par leurs parents,
    \end{itemize}
    \begin{center}\rule{0.7\textwidth}{0.2pt}\end{center}
    On peut alors dire que :
    \begin{itemize}[label=\textbullet]
        \item la population statistique est \answer{les élèves}
        \item un individu est \answer{un élève}
        \item la variable (le caractère) est \answer{le moyen de transport utilisé}
        \item l'effectif des individus venant à pied est \answer{125}
    \end{itemize}
\end{exemplebox}

\begin{remarquebox}
    La représentation brute des résultats n'est guère exploitable, il est 
    donc usuel de regrouper les résultats par valeurs identiques dans un tableau
    appelé \hlb{tableau des effectifs}.
\end{remarquebox}

\begin{exempleboxsuite}
    On peut reprendre les valeurs de l'exemple 1 dans un tableau des
    effectifs: 

    \vspace{5pt}

    \begin{tabular}{|c|c|c|c|c|}
        \hline
        \bb{Transport} &
        2 roues & Bus & A pieds & Parents \\
        \hline
        \bb{Nombre d'élèves (= Effectif)} &
        50 & 250 & 125 & 75 \\
        \hline
    \end{tabular}
\end{exempleboxsuite}


